\documentclass[11pt,a4paper]{article}
\usepackage[margin=2.4cm]{geometry}
\usepackage{booktabs,array,tabularx}
\newcolumntype{Y}{>{\raggedright\arraybackslash}X}
\usepackage{amsmath,amssymb}
\usepackage[hidelinks]{hyperref}
\usepackage{xcolor}
\usepackage{enumitem}
\usepackage{caption}
\usepackage{float}
\captionsetup[table]{position=top,skip=4pt,font=small}
\emergencystretch=1.5em

% ---------------------------------------------------------------------------
% Freeze switch. Release candidate: \frozenfalse. Deposit: \frozentrue, then
% fill every \todo (grep -n 'todo{' must return nothing) and rebuild.
\newif\iffrozen \frozenfalse
\iffrozen
  \newcommand{\prop}[1]{#1}
  \newcommand{\todo}[1]{\errmessage{Unfilled placeholder: #1}}
\else
  \newcommand{\prop}[1]{\textcolor{blue}{#1}}
  \newcommand{\todo}[1]{\textcolor{red}{[TBD: #1]}}
\fi
\newcommand{\claim}[1]{\textbf{#1}}
% ---------------------------------------------------------------------------

\title{Pre-registered protocol: a global sensitivity audit of the
scenario-generation pipeline of Li et al.\ (2025), \emph{Nature Climate
Change}\\[4pt]\large Modelling of the Modelling Process (MoMP) series:
optimal-transport GSA case}
% Author block kept out, as in the version circulated on 21 September 2026.
% To include it, replace \author{} with the block below and confirm the order.
% \author{Samuele Lo Piano$^{1,2}$ \and Giulia Vannucci$^{3}$ \and Arnald Puy$^{4}$
%   \and Roberta Siciliano$^{3}$ \and Andrea Saltelli$^{2,5}$\\[6pt]
%   \parbox{0.92\textwidth}{\centering\small
%   $^{1}$Faculty of Management and Economics, Gda\'nsk University of Technology, Poland\\
%   $^{2}$Barcelona School of Management, Universitat Pompeu Fabra, Spain\\
%   $^{3}$Department of Electrical Engineering and Information Technology,
%         University of Naples Federico II, Italy\\
%   $^{4}$School of Geography, Earth and Environmental Sciences, University of Birmingham, UK\\
%   $^{5}$Centre for the Study of the Sciences and the Humanities, University of Bergen, Norway}}
\author{}
\date{Version \iffrozen 1.0\else 0.2 (release candidate)\fi, \todo{date of deposit},
DOI: \todo{reserved on Zenodo before upload}}

\begin{document}
\maketitle

\iffrozen\else
\noindent\fbox{\parbox{\dimexpr\textwidth-2\fboxsep-2\fboxrule}{\small
Release candidate. Values agreed in the co-authors' review of 23 October 2026
are in black. Text in \prop{blue} was added while transcribing that review and
awaits approval. Items in \todo{} are filled from the Phase~0C adjudication or
at deposit. This box, all colour and all placeholders disappear in the frozen
version.}}
\fi

\begin{abstract}
\noindent We pre-register a global sensitivity analysis (GSA) of the
scenario-generation pipeline of Li, Zhu, McJeon, Byers, Zhou and Ou (2025),
``Using deep learning to generate key variables in global mitigation
scenarios'', \emph{Nat.\ Clim.\ Change} 15:760--768
(\path{doi:10.1038/s41558-025-02352-8}). The object of analysis is not the
climate or the AR6 ensemble but the pipeline itself. The analyst's
methodological choices are the input factors, the generated multivariate
ensemble and the paper's validation claims are the outputs, and sensitivity
is measured by optimal-transport (Wasserstein) displacement of the synthetic
ensemble, estimated from given data. Before any GSA run, this protocol
freezes the operational rules for the paper's claims
(\S\ref{sec:claims}), the factor space (\S\ref{sec:factors}), the
experimental design and compute budget (\S\ref{sec:design}), the estimators
and inference procedure (\S\ref{sec:estimators}), and the decision criteria
(\S\ref{sec:criteria}). The completed replication record (Phase~0,
\S\ref{sec:phase0}) is reported for transparency. At the time of this
deposit, no claim rule (X1a--X5) has been evaluated at any factor setting,
including the paper's own values, no retraining has been performed, and no
noise-floor replicate has been run. The only pipeline stage executed is the
deterministic regeneration check of Phase~0B, which evaluates no claim.
\end{abstract}

\section{Purpose and scope}
Continuing the MoMP programme, this study asks whether the properties Li
et al.\ attribute to their synthetic-scenario framework are properties of
the framework or of the specific, partly undocumented methodological
choices through which it was instantiated. We audit the epistemic
apparatus, not the climate--economy content: factors are analyst choices,
outputs are (i) the synthetic ensemble as a multivariate distribution and
(ii) pass/fail evaluations of the paper's claims under fixed operational
rules.

\section{Audited artefacts}\label{sec:artefacts}
All artefacts are pinned by version and SHA-256 in \texttt{DATA.md} of the
project repository, archived with this deposit. The principal items follow.
\begin{description}[nosep,font=\normalfont\bfseries,leftmargin=3.8em,labelwidth=3.4em,labelsep=0.4em,align=left]
\item[A1] Deep-IAM code, v1.0.0, \path{doi:10.5281/zenodo.15280653}, the
  version cited by the paper and the audited object.
\item[A2] Deep-IAM v1.0.1 (2026-04-25), admitted as a documented supplement
  only for its two added files: \texttt{Secondary-data-generate.ipynb} (the
  generator of the deposited synthetic data, absent from v1.0.0) and
  \texttt{Configuration of Environments.txt} (environment pins: Python
  3.8.20, TensorFlow 2.13.0, tsgm 0.0.7, \prop{TensorFlow Probability
  0.21.0,} scikit-learn 1.3.2, XGBoost 2.1.1). Full diff archived.
\item[A3] Synthetic-scenario deposit, \path{doi:10.5281/zenodo.15240553}
  (Gen and Real CSVs, validated in Phase~0).
\item[B1--B2] AR6 Scenarios Database v1.1 (metadata and World
  timeseries), \path{doi:10.5281/zenodo.7197970}.
\end{description}

\section{Phase 0: completed replication record}\label{sec:phase0}
Completed before this deposit. Machine-readable reports are archived under
\texttt{results/}.
\begin{enumerate}[nosep]
\item \textbf{Metadata fidelity.} The AR6 v1.1 vetted-with-climate
  population reproduces the paper's Extended Data Table~1 exactly: $n=1202$,
  with per-category counts matching digit for digit
  (\path{phase0_metadata_check.json}).
\item \textbf{First undeclared filter reconstructed (D2).} The training
  population of SI~S3.1 (680/308/172, $n=1160$) is reproduced exactly by
  complete 2020--2100 coverage on the five feature variables
  (\path{Primary Energy|Coal}, \path{|Oil} and \path{|Gas},
  \path{Final Energy|Liquids}, \path{Carbon Sequestration|CCS}), with Kyoto
  Gases exempt. This rule appears nowhere in the paper or SI.
\item \textbf{Second undeclared filter reconstructed (D9).} The population
  underlying the paper's figures and label-transfer benchmark (Fig.~4a,
  ED~Fig.~2, Table~1: 554/237/110, $n=901$) is reproduced exactly by
  complete 2020--2100 coverage on the five feature variables and all ten
  electricity fuel-level variables, Kyoto Gases exempt, raw
  (pre-imputation) CS availability required. Across the two silent drops
  the smallest category (C78) loses 43\% of its vetted scenarios
  (193$\to$110).
\item \textbf{Deposit validated.} After harmonizing a Mt-versus-Gt unit
  inconsistency (deposit in Mt, the paper's Source Data Fig.~3 in Gt), all
  1{,}620 quantile comparisons (6 variables $\times$ 6 series $\times$ 9
  years $\times$ 5 quantiles) pass at 2\% relative tolerance. Gen files
  hold 10{,}000 scenarios per category.
\item \textbf{Version provenance (D7).} v1.0.0 cannot reproduce the
  deposited synthetic data or Fig.~5. The generating notebook appeared in
  v1.0.1 ten months after publication, without changelog.
\item \textbf{Non-reproducibility of training (D1).} The authors state that
  TSGM training is not reproducible even under fixed seeds and distribute
  trained weights, which motivates the noise-floor design of
  \S\ref{sec:design}.
\item \textbf{Phase 0B, executability and determinism of the generation
  stage.} Evaluating no claim rule, we executed the pipeline's one
  nominally deterministic link. The authors' shipped decoder weights
  (\texttt{vae\_conv5}, 9 steps $\times$ 16 features, latent dimension 10)
  were loaded, the \texttt{TSFeatureWiseScaler} was refit on the deposited
  Real CSVs (which the D9 reconstruction identifies as the training
  tensors), the latent noise was drawn under the documented global seed 3,
  decoded, and compared with the deposited Gen CSVs (environment: Python
  3.10.21, TensorFlow 2.13.0, tsgm 0.0.7, macOS/arm64, whereas the authors
  used Windows x86-64 and Python 3.8; report
  \path{phase0b_regeneration.json}). Executability is confirmed: the weights
  load, the architectures match and generation runs end to end from the
  archived artefacts alone. Sample-level reproduction fails under all three
  latent-stream readings of the documentation (TensorFlow global seed, a
  second-execution probe and the SI's NumPy reading), with median
  elementwise relative deviation 0.33--0.45, i.e.\ entirely different
  draws. The distributions agree in the bulk, with median quantile
  deviation 1.8\% of the reference spread (Q90--Q10) and 95th percentile
  10.3\%, yet about 4.4 times the Monte Carlo floor obtained from two
  independent 10{,}000-draw ensembles through the same decoder, the excess
  being concentrated in distribution tails and tight-spread cells (e.g.\
  carbon sequestration in 2020). A scaler-bounds diagnostic flags no
  feature, consistent with the Real CSVs being the scaler-fit data. The
  residual is most plausibly a checkpoint mismatch (the shipped weights
  differing from those that generated the deposit), consistent with the
  repository README's account of multiple training sessions with
  ``comparable'' results. The generation stage is therefore executable and
  distributionally faithful, but the deposited outputs cannot be tied to
  the documented seed or, apparently, to the shipped weights. The
  determinism gap the authors acknowledge for training (D1) extends to
  generation (D11). This check answers no question posed in
  \S\ref{sec:claims}--\S\ref{sec:criteria}.
\item \prop{\textbf{Phase 0C, static adjudication from code and real data.}
  Without executing any model, the v1.0.0 notebooks are scanned for the GAN
  training budgets (D4), the configuration behind the GAN baselines of
  Table~1 (D12), the label-transfer classifier settings (G8) and the
  sum-check aggregation (X3), and the AR6 reference value of the
  electricity sum-check is recomputed from the deposited Real CSVs alone
  under candidate aggregations, to identify the one that reproduces the
  paper's 1.2\%. No synthetic file is read in this step. Outcome:}
  \todo{transcribe \path{phase0c_code_scan.txt} and
  \path{phase0c_sumcheck_reference.json}}.
\end{enumerate}
The full inconsistency catalogue (D1--\prop{D12}) is archived as
\texttt{DISCREPANCIES.md}.

\section{Claims and operational rules}\label{sec:claims}
Each claim of the target paper is translated into a rule verifiable on the
pipeline's output. Populations are declared explicitly, given D2 and D9.
\prop{Claims are evaluated on the three aggregated categories (C1234, C56,
C78) and on the evaluation set $E$ of \S\ref{sec:estimators}. When factor G2
uses finer categories, the synthetic ensembles are pooled to the three
aggregates with weights equal to the real category counts in the 901-set,
and k-means clusters are mapped to the aggregate they overlap most.} No rule
will be altered after deposit except under the deviations policy
(\S\ref{sec:deviations}).

\begin{description}[style=nextline]
\item[\claim{X1a} (distributional match).] ``[The VAE] closely matches the
  original AR6 scenarios in terms of mean and the spread, particularly for
  C1234.'' Rule: for each of the six key variables and each category,
  \prop{the median, over the nine decades and the five quantiles (Q10, Q25,
  Q50, Q75, Q90), of the absolute difference between synthetic and real
  quantiles divided by the real 10--90 percentile range of that variable
  and decade (the band drawn in the paper's Fig.~3)} is at most
  $\tau_{1}=0.15$, the real population being the 901-set. ``Particularly
  C1234'' holds if the C1234 deviation, \prop{pooled over variables,
  decades and quantiles,} is the smallest of the three categories.
  \prop{This normalization replaces the deviation relative to the quantile
  itself, which is ill-conditioned where quantiles approach or cross zero
  (coal in C1234 late in the century, Kyoto gases in C1234 around net zero,
  carbon sequestration in C78). In the Phase~0B check, a
  synthetic-versus-synthetic comparison, relative deviations reached about
  1{,}260\% (electricity from coal, C1234, 2070, Q10) in a category where no
  spread-normalized deviation exceeded 0.34.}
\item[\claim{X1b} (narrower distributions).] ``Synthetic scenarios have
  narrower distributions for all climate categories.'' Rule:
  IQR$_{\mathrm{syn}} <$ IQR$_{\mathrm{real}}$ for all six key variables and
  all three categories in 2050 and 2100.
\item[\claim{X2} (correlation replication).] ``The synthetic scenarios also
  retained the correlations between key variables \dots\ matching closely
  those observed in the original AR6 scenarios.'' Rule, for C1234, the
  category shown in the paper's Fig.~4: (i) the maximum absolute
  elementwise difference between the synthetic and real correlation
  matrices of the six key variables, taken over the 15 pairs and the three
  years 2030, 2050 and 2100 (45 elements), is at most $\tau_{2}=0.15$, and
  (ii) the paper's reported temporal ordering (coal--Kyoto correlation
  stronger in 2030 than in 2100) holds in the synthetic data. The sampling
  noise of a correlation is at most about $1/\sqrt{n}$, i.e.\ about 0.04 at
  $n=554$, hence about 0.13 on the maximum of 45 elements, and larger in
  practice because scenarios from the same IAM are not independent.
  $\tau_{2}$ is a heuristic ceiling, since the real 901-set is the training
  data and the generator fits that sample.
\item[\claim{X3} (internal consistency).] ``Among the 30{,}000 synthetic
  scenarios, 258 exhibited lower sum-check errors than the 1.2\% average
  error in the AR6 Scenarios Database.'' Rule: \prop{the fraction} of
  synthetic scenarios with electricity sum-check error below 1.2\% is at
  least \prop{$258/30{,}000 = 0.86\%$}. The sum-check error of a scenario
  is $|\mathrm{Gen}_{\mathrm{total}}-\sum_i
  \mathrm{Gen}_i|/\mathrm{Gen}_{\mathrm{total}}$ (the paper's Eq.~3),
  aggregated over the nine decades by \todo{the mean, unless Phase~0C
  identifies the authors' aggregation}. \prop{Stating the benchmark as a
  fraction prevents the sample-size factor (G7) from moving the rule
  mechanically.}
\item[\claim{X4} (VAE superiority).] ``VAE achieved higher category label
  transferring accuracy \dots\ compared to RCGAN and CGAN'' (Experiment~I,
  synthetic$\to$real, test population the 901-set). Rule: accuracy(VAE)
  $>$ accuracy(RCGAN) and accuracy(VAE) $>$ accuracy(CGAN) in the same
  run. \prop{The three generators are trained in the same run on the same
  population and the same variables, so the comparison is like-for-like by
  construction, whatever configuration the published baselines used
  (D12).}
\item[\claim{X5} (bias inheritance, secondary).] The paper positions
  synthetic generation as mitigating ``model-determined biases'' and
  ``non-random submission''. Rule: a source-model classifier (nine model
  groups, merged as in Dekker et al.) is trained on the real 901-set and
  applied to the synthetic scenarios. X5 (inheritance) holds if, per
  category, the total-variation distance between the predicted model mix
  of the synthetic ensemble and the real ensemble's model shares is smaller
  than the distance to the uniform mix. A baseline model-identifiability
  check on real data (held-out accuracy above the majority-class rate)
  gates the test.
\end{description}

\section{Factor space}\label{sec:factors}
Factors, paper values and perturbation ranges are frozen in
\path{config/paper_values.yaml}, archived with this deposit and summarized
in Table~\ref{tab:factors}. Expensive factors require retraining, cheap
factors vary per trained model.

\begin{table}[H]
\caption{Factor space: paper values with provenance, ranges and cost.}
\label{tab:factors}
\footnotesize
\begin{tabularx}{\textwidth}{@{}l Y Y l@{}}
\toprule
Factor & Paper value (provenance) & Range or levels & Cost\\
\midrule
G1 population & \prop{901-rule (D9, undeclared; see D12)} &
  \prop{\{901; 910, CS imputed; 828, Kyoto required; 1{,}160, D2 rule;
  1{,}202, all vetted\}, missing values filled by G4} & exp.\\
G2 aggregation & C1234/C56/C78 (Methods) &
  \{3-way; 4-way; 8-way; k-means labels\}\prop{, pooled for evaluation
  (\S\ref{sec:claims})} & exp.\\
G3 variable selection & 5 key features ranked by RF (SI S2) &
  pool completeness threshold 0.6--0.9; number of key features 4--9\prop{,
  generated set united with $E$} & exp.\\
G4 preprocessing & mean imputation by year $\times$ category; (0,1)
  scaling & imputation \{mean yr$\times$cat.; mean yr; median
  yr$\times$cat.; none\}; scaling \{(0,1); standardized\} & exp.\\
G5 architecture & VAE latent dimension 10; ED Table~3 & latent dimension
  4--20; epochs $\times$ 0.25--2 & exp.\\
G6 seed & seed 3; training non-reproducible (D1) & \prop{integer seed,
  uniform; dummy reference (\S\ref{sec:criteria})} & exp.\\
G7 sample size & 10{,}000 per category & 1{,}000--30{,}000 per category &
  cheap\\
G8 transfer classifier & RF (code) & \{RF as in code; RF defaults;
  logistic; XGBoost\} & cheap\\
\bottomrule
\end{tabularx}
\end{table}

\noindent \prop{Factor G9 (validation metrics) of the circulated version is
withdrawn: its range was never specified, and varying the thresholds of
the claim rules would turn the pre-registered rules into moving targets.
The thresholds are frozen in \S\ref{sec:claims}, and their influence is
reported descriptively (\S\ref{sec:estimators}).} The D4 documentation
conflict (CGAN and RCGAN epochs, SI versus ED Table~3) is handled as an
explicit two-level factor nested in G5 unless adjudicated \prop{from the
code in Phase~0C} before the freeze.

\section{Design and compute budget}\label{sec:design}
\begin{enumerate}[nosep]
\item \textbf{Noise floor.} $R_{0}=20$ seed-only replications at paper
  values for all factors. The floor $\delta_{0}$ is the 95th percentile of
  pairwise distances between replicate ensembles, per category. \prop{The
  replicates also give the pass rate of every claim at the paper's own
  configuration, i.e.\ the statistical replication of the published
  results.}
\item \textbf{Outer design.} \prop{Scrambled} Sobol' quasi-random sample of
  size $N_{\mathrm{outer}}=512$ over the expensive factors G1--G6, one
  training per point, environment pinned to A2's configuration.
\item \textbf{Inner design.} \prop{Per trained model,
  $N_{\mathrm{inner}}=8$ points of an independently scrambled Sobol'
  sequence over the cheap factors G7 and G8. Each trained generator
  produces 30{,}000 scenarios per category once, subsampled to the G7 value
  of each inner point.}
\item \textbf{Stopping rule.} The budget is
  $N_{\mathrm{outer}}\times$(training cost). If bootstrap CIs on the
  leading OT indices have half-width below 0.05 at $N_{\mathrm{outer}}/2$,
  the design may stop early. It may not be extended beyond
  $N_{\mathrm{outer}}$ without a documented deviation.
\item \textbf{Outputs stored per run:} \prop{per category, the standardized
  quantile vectors on $E$, the 200 projected quantile vectors and a
  2{,}000-point subsample of the 10{,}000-scenario ensemble}, all claim
  evaluations X1a--X5 \prop{for each inner point}, wall-clock time, and the
  full factor vector.
\end{enumerate}

\section{Estimators}\label{sec:estimators}
\paragraph{Ensemble-level sensitivity.} Each training yields, per
aggregated category, an ensemble of 10{,}000 synthetic scenarios
\prop{(G7's paper value) represented on the evaluation set $E$: the six key
variables and the ten electricity variables over the nine decades
2020--2100.} Each variable--year marginal is standardized by its standard
deviation in \prop{the corresponding category of} the real 901-set. The
distance between ensembles is the exact one-dimensional Wasserstein
distance on each variable--year marginal, aggregated by the mean \prop{of
squared Wasserstein-2 distances and square-rooted, so that it is the
Euclidean distance between the vectors of standardized quantile functions
(99 probability levels)}. It is complemented by the sliced Wasserstein
distance on the joint, with $L=200$ projection directions fixed once,
shared across all runs and replicates \prop{and archived with the
harness}, so that projection noise is not confounded with training noise.
\prop{Both distances embed the output $Y$ of each training as a point of a
Euclidean space, and the sensitivity index of factor $X_j$ is the
given-data optimal-transport index of the Borgonovo--Chiani--Plischke class
with quadratic cost,
\[
\iota_j \;=\; \frac{\mathbb{E}\big[W_2^2\big(P_Y,\,P_{Y\mid X_j}\big)\big]}
{\mathbb{E}\,\lVert Y-Y'\rVert^2},
\]
which lies in $[0,1]$, vanishes under independence and equals one when $Y$
is a function of $X_j$. It is a first-order (main-effect) index.}
Conditional laws are estimated with $M=8$ bins at $N=512$ (about
$N^{1/3}$, 64 runs per bin), a heuristic choice, with results also
reported for $M\in\{5,8,12\}$. Discrete factors use exact level
partitions, with the partition-collapse guard of our prior OT-GSA
implementation. \prop{Indices are computed per aggregated category and for
the three categories jointly.} As a robustness check, the \prop{debiased}
Sinkhorn divergence is computed with $\varepsilon = 0.05$ times the median
pairwise cost (squared Euclidean distance in the standardized space),
computed on subsamples of 2{,}000 points per ensemble. \prop{$\varepsilon$ is
computed once, on the deposited synthetic ensembles (A3), and held fixed
for all comparisons. The index is then obtained with this divergence as
the ground cost between runs.}

\paragraph{Claim-level sensitivity.} For each rule X1a--X5, the fraction of
runs in which the claim holds, \prop{over the full design and, separately,
over the $R_0$ replicates at the paper's values. Given-data first-order
indices $S_j$ of the binary pass indicator are estimated with the same
partitions, as the between-bin variance of the pass rate over its total
variance. Total-order indices are not estimable from a given-data design at
this budget and are not reported. The interaction share $1-\sum_j S_j$ is
reported in their place, and the $S_j$ remain comparable to the first-order
indices of the companion analysis of Li P.-H.\ et al.\ (2023).}

\paragraph{Threshold sensitivity.} \prop{The pass fraction of each claim is
also reported as a function of its threshold ($\tau_1$ and $\tau_2$ from
0.05 to 0.30, sum-check threshold from 0.8\% to 2.0\%). These curves are
descriptive and enter no decision.}

\paragraph{Inference.} Bootstrap with $R=1000$ resamples, with 95\%
intervals on all reported indices. \prop{The resampling unit is the
training (outer point), each resampled training carrying its eight inner
points, so that evaluations sharing a trained generator are not treated as
independent.}

\section{Decision criteria}\label{sec:criteria}
\prop{A factor is influential for the ensemble output if the lower bound of
the 95\% interval of its index exceeds the upper bound of the 95\% interval
of the index of the training seed (G6). Since training is not reproducible
even under a fixed seed (D1), the seed acts as a dummy factor: its index
measures what training stochasticity and finite-sample estimation produce
under the same estimator and partition. The same criterion applies to the
claim-level indices. $\delta_0$ is reported as the distance below which two
ensembles cannot be told apart from retrained replicates.} A claim is
robust if it holds in at least 80\% of runs across the design, and fragile
if it holds in fewer than 50\%. The study's headline question, whether the
paper's findings are properties of the framework or of the analyst's
choices, is answered by the joint pattern of (i) claim fragility and (ii)
the identity of the dominant factors (declared versus undeclared choices).

\section{Deviations policy}\label{sec:deviations}
Any departure from this protocol is logged in \texttt{DEVIATIONS.md} with
date, reason, and the analysis state at the time of the change. Results
affected by deviations are reported separately from the pre-registered
analysis. Adjudications of placeholders completed before the Zenodo freeze
are part of this protocol. The deposited PDF contains no placeholders.

\section{Archive contents at deposit}
Protocol PDF, \texttt{DATA.md}, \texttt{DISCREPANCIES.md},
\prop{\texttt{DEVIATIONS.md} (empty),} \path{config/paper_values.yaml}
(frozen), Phase-0 reports (\texttt{results/*.json}, including the
Phase~0B \prop{and Phase~0C} reports), v1.0.0/v1.0.1 diff, commit hash of
the analysis harness (\todo{hash}). GSA results: none, no claim rule
evaluated at any factor setting (the paper's own values included), no
retraining performed, no noise-floor replicate run.

\end{document}
